% Section 6: monotone EI, N1, N2 (PHASE2c) and PACKAGE_DESIGN section 5
\section{Monotone ecological inference: signs, N1 and N2}\label{sec:monotone}

Source: PHASE2c (Appendix~\ref{app:phase2c}), following Elzayn, Goldin, Guage,
Ho and Morton~\citep{egghm2025} (EGGHM, arXiv v1 numbering), and PACKAGE\_DESIGN
\S\S2.2 and 5 (Appendix~\ref{app:pkgdesign}). Checks: \texttt{verify\_phase2c.py},
29/29 PASS, each Batch~B correction tagged [F], [V] or [I]. Two settings are used:
\begin{itemize}[nosep]
\item \textbf{Finite population (N1).} The design is fixed; unit weights
$p_i=N_i/\sum_jN_j$; the realised rates satisfy $\beta_i\in S(X_i,T_i)$; the
target is the actual rate $B$. All N1 results read the unit tables as
\emph{continuous} (real cell masses).
\item \textbf{Superpopulation (K-rows and N2).} $m(v)=\E[\beta\mid V=v]$; the
target is $\tilde B$; a superscript $^*$ marks an association computed with $m$.
\end{itemize}

\subsection{The associations}

Fix a column and write $w^r_i=p_iX_{ir}/\mu_r$.
\begin{itemize}[nosep]
\item \emph{Between, pairwise:} $\delta_B(r,s)=E_p[X_rX_s(\beta_r-\beta_s)]$;
\emph{one-vs-rest:} $\delta^r_B=\sum_{s\ne r}\delta_B(r,s)$.
\item \emph{Within, W2 (pooled):} $\Delta=\sum_s\mu_s\Cov_{w^s}(X,\beta_s)$, one
sum of row-specific covariances.
\item \emph{Within, W1 (one-vs-rest):}
$\delta^r_W=\mu_r\Cov_{w^r}(\beta_r,X_r)+(1-\mu_r)\Cov_{w^{\neg r}}(\beta^{\neg r},X_r)$,
with A's share as the predictor in both terms.
\end{itemize}
\begin{plain}
The \emph{between} association says that, on average within the same unit, group
A is more likely than the others to choose the outcome. It is a weighted average
with weights $p_iX_{ir}(1-X_{ir})$: pure units get no weight and a large mixed
unit can dominate. It is not an ordering in every unit or against each rival. The
\emph{within} association says that rates rise, on average, with A's share. W1 and
W2 can have opposite signs while every group's rate rises (EI referee's example:
$\delta^A_W=-0.02633$, $\Delta_A=+0.00110$, check C14). A sign on W2 signs one
pooled sum, not each group. Covariances also reflect sorting; a causal reading
needs a separate design.
\end{plain}
\begin{revnote}
(P-E1--P-E5, P-E8, P-B28.) W2 was earlier described as ``each group
separately''; the weighting of $\delta_B$, the predictor in W1, sorting, the
Goodman-bias relation and the status of W2, pair and CR as extensions of EGGHM
were added as remarks on Definition~1.
\end{revnote}

\begin{result}[Lemma M1, identities --- PHASE2c \S1]\label{res:m1}
\status{Proved.} Under A1 only (with the stated per-part conditions):
(a) $\mathbf 1^\top\Delta=0$ and Goodman's $g$ satisfies $g-B_{\cdot c}=M^{-1}\Delta$;
(b) $NM_r-B_r=-\delta^r_B/\mu_r$;
(c) the one-vs-rest collapse is a valid two-group population, so
$ER^r-D^r=\delta^r_W/\Var_p(X_r)$ and $(NM_r-NM_{\neg r})-D^r=-\delta^r_B/(\mu_r(1-\mu_r))$;
(d) $B$, every $\delta_B$, $\Delta$ and $\delta^r_W$ are affine in the unit rates.
\end{result}
\begin{plain}
These are accounting identities. The neighbourhood model (NM) misses the true
gap exactly by the between association; the aggregate regression slope (ER) misses
it exactly by the within association. So a sign belief about either association
is a one-sided bound relative to a number computed from the data.
\end{plain}
\proofref{Appendix~\ref{app:phase2c}, \S1, proofs of (a)--(d); check C1}
\begin{revnote}(P-B28.) The invertibility of $M$ is needed for (a) only.\end{revnote}

\subsection{N1: sign restrictions on realised tables}

\begin{result}[Proposition M2, N1 is sharp for continuous unit tables --- PHASE2c \S2]\label{res:m2}
\status{Proved} (continuous reading). For $F=\{\beta:\beta_i\in S(X_i,T_i)\cap
U_i,\ g_k(\beta)\ge-\varepsilon_k\}$: (a) the values of $\varphi$ attained by
continuous unit tables are $\varphi(F)=[\min_F\varphi,\max_F\varphi]$, sharp
given the margins; (b) $F=\emptyset$ iff the restrictions contradict the margins
in the continuous problem; (c) under CR the set is a union of at most two
intervals; (d) the sets are nested componentwise in $\varepsilon$ with continuous
piecewise-linear endpoints; the preserving set of slacks is in general not a
product of intervals, and its boundary is the breakdown frontier; a scalar
breakdown value needs a stated path.
\end{result}
\begin{plain}
With sign beliefs alone and no regression, the bound is an LP over all unit
tables that fit the margins, and it is exactly right for real-valued tables.
With whole-number ballots it can differ (Remark~M2.1).
\end{plain}
\proofref{Appendix~\ref{app:phase2c}, \S2, Proposition~M2 steps 1--6}
\begin{revnote}
(P-B1, P-B2, I.10.) ``Sharp in the finite-population sense'' became ``sharp for
continuous unit tables, conditional on the margins''; ``the breakdown value'' for
vector slacks became the breakdown frontier with a scalar value along a path; in
round~2, the set and its frontier were distinguished. Checks C7 [F], C7b [V].
\end{revnote}

\begin{result}[Remark M2.1, integer counts --- PHASE2c \S2]\label{res:m21}
\status{Proved} for $2\times2$ with margins only; \status{open} otherwise (R15).
In $2\times2$ with margins only and sign rows on one column, the integer gap set
is the DD interval intersected with a lattice of step $1/N_1+1/N_0=1/(N\mu(1-\mu))$;
endpoints differ from the continuous ones by less than one step (e.g.\
$4.8\times10^{-6}$ for $N=10^6$, $\mu=0.3$). With unit restrictions, joint
columns or $R\ge3$, the integer--continuous gap is not bounded by this argument.
\end{result}
\begin{plain}
Counting whole ballots matters only for tiny tables or near-ties, at least in
$2\times2$; the referees' four-voter examples show that CR can then be empty and
an endpoint can move.
\end{plain}
\proofref{Appendix~\ref{app:phase2c}, \S2, Remark~M2.1 (proposition and proof); checks C7, C7a}
\begin{revnote}(P-B1, P-E11, I.5.) Added after review; the lattice step quoted with both group totals in round~2.\end{revnote}

\begin{result}[Proposition M3, two rows --- PHASE2c \S2]\label{res:m3}
\status{Proved.} For $R=2$, $\Var_p(X)>0$, continuous tables, no unit
restrictions and no other sign rows: (a) $\delta_W\ge0\iff D\le D_{ER}$ and
$\delta_B\ge0\iff D\ge D_{NM}$; (b) $\{\delta_W\ge0\}$ gives
$[D^-,\min\{D^+,D_{ER}\}]$, empty iff $D_{ER}<D^-$; (c) $D_{NM}\in[D^-,D^+]$
always, so between signs are never refuted in this problem; (d)
$D_{NM}=\gamma D_{ER}$, $\gamma\in(0,1]$, and CR has exactly one non-empty branch
when $\gamma<1$ and $D_{ER}\ne0$; at $\gamma=1$ or $D_{ER}=0$ both branches are
non-empty with the same gap image. The CR set is $[D_{NM},\min\{D^+,D_{ER}\}]$
if $D_{ER}\ge0$ and $[\max\{D^-,D_{ER}\},D_{NM}]$ otherwise; under it
$\operatorname{sign}D=\operatorname{sign}D_{ER}$.
\end{result}
\begin{plain}
In this base problem (two groups, real-valued tables, nothing else imposed) only a
within sign can be contradicted by the margins; a between sign and contextual
reinforcement are always feasible, because ``every group votes like its unit''
fits them. That non-refutability belongs to this base system only: under the
linear rung CR can be rejected jointly (Remark~R3.2); with integer counts CR can be
rejected outright; known cells or unit orderings remove the witness. The
equal-rates table satisfies the CR \emph{union}, in R$\times$C too, but only the
branch matching its within sign.
\end{plain}
\proofref{Appendix~\ref{app:phase2c}, \S2, Proposition~M3 steps 1--8 and ``Reading''; checks C2, C8, C8b, C15}
\begin{revnote}
(P-B3--P-B5, P-E12, P-E13; round~2 I.1, I.9.) The hypotheses ``no unit-level
restrictions, no other sign rows'' were missing; (d) failed at $\gamma=1$;
``CR is never refuted'' now names its base system; branches and their gap images
are distinguished; in round~2 the equal-rates witness was corrected to give the
CR union, not both branches (three-row example, C15 [F]).
\end{revnote}

\begin{result}[Proposition M4, joint restrictions --- PHASE2c \S2]
\status{Proved.} Imposing signs on several columns or rows jointly gives a subset
of each single restriction's feasible set.
\end{result}
\begin{plain}
Joint beliefs are weakly tighter. In $2\times3$ simulations, signing two columns
jointly narrowed column 1 strictly in 75 of 139 populations. With two
complementary columns, the \emph{same} weak within sign on both forces $D=ER$.
\end{plain}
\proofref{Appendix~\ref{app:phase2c}, \S2, Proposition~M4; checks C3, C3b}
\begin{revnote}(I.8.) The round-1 text said ``opposite within signs''; it is the same sign on complementary columns.\end{revnote}

\subsection{Signs under linear contextual effects (K-rows)}

\begin{result}[Proposition M5 --- PHASE2c \S3]\label{res:m5}
\status{Proved}; (c) \status{for M0 only}. Write the rung's means as
$m_i=m_{\rm id}(V_i)+L_ik$. (a) Every association $\delta^*$ is affine in $k$ under
every rung, so each sign is one LP row and CR is a union of two LPs. (b) Mean-level
ordering in every district on its support implies $\delta^*_B(r,s)\ge0$, pooled.
(c) For $R=2$ under M0, $\delta^*_W=a(w_1+d_1)+bw_1$ with $a,b\ge0$, so both
slopes $\ge0$ imply $\delta^*_W\ge0$. Under district rungs the pooled
$\delta^*_W$ contains a between-district term with no sign, and the implication
fails (six-unit example: within districts $+0.00064$, pooled $-0.04927$).
\end{result}
\begin{plain}
Under a linear rung the sign beliefs become cheap extra constraints on the
unknown $K$. But a pooled sign under district models is one restriction across
all districts: positive slopes in every district do not make the pooled within
association positive, because district baselines can fall as composition rises.
\end{plain}
\proofref{Appendix~\ref{app:phase2c}, \S3, Proposition~M5 steps 1--5; checks C4 (rebuilt), C9}
\begin{revnote}
(P-B7; E2-4.) M5 was stated for every rung with the pooled formula. Now (a) uses
the stacked parameterisation of each rung, (b) needs district-specific ordering,
(c) holds for M0 only. Sharpness and inference for M-SIGN under district rungs
became proof item R16 (open). A district rung with sign rows has no coverage
theorem yet.
\end{revnote}

\paragraph{Remark M5.1 (which $\delta$).} N1 uses the realised associations;
K-rows use $\delta^*$. Their difference is a design-weighted sum of unit residuals
with mean zero; under P3-S1a its standard deviation is at most
$\tfrac12\{\sum_i(\sum_r|c_{ir}|)^2\}^{1/2}$. The earlier ``order $p^{-1/2}$''
claim needed conditions (P-B8).

\subsection{N2: monotone nonparametric contextual effects (two rows)}

N2 is defined on the finite empirical design law (support points
$x_1<\dots<x_{J_X}$). Model N2($\sigma$): (S1), $\E[\beta\mid V]=m(X)$, and for
designated $(r,c)$, $\sigma_{rc}m_{rc}$ is nondecreasing on the support.

\begin{result}[Lemma N2.1 and Proposition N2.2 --- PHASE2c \S4]\label{res:n22}
\status{Proved} (on the finite support). N2.1: $m(x)\in S(x,\tau(x))$ at every
support point. N2.2: (a) the LP over $\mathcal M_O$ (support-point means in
$S(x_j,\tau(x_j))$ plus adjacent order rows) is an \emph{outer} set for $\tilde
B$, not claimed sharp; (b) every point of $\mathcal M_O$ extends to a monotone
Borel $m$, so on the complete finite support nothing is lost (on a coarser grid
it is only a relaxation); (c) with quadratic $\tau$, M0 with signed slopes lies
inside N2; (d) both rows nondecreasing imply $\delta^*_W\ge0$; (e) for $R\ge3$ a
direction needs a simplex chart, and (d) has no general analogue.
\end{result}
\begin{plain}
N2 assumes only ``no reversals'': each group's average rate moves in one
direction as its share grows, in any shape. It is weaker than straight lines with
signed slopes, so its bound contains theirs. It is not in the package.
\end{plain}
\proofref{Appendix~\ref{app:phase2c}, \S4, Lemma~N2.1 steps 1--2; Proposition~N2.2 steps 1--5; checks C5, C10, C10b}
\begin{revnote}
(P-B9, P-B10, P-B29.) The support convention was added; (b) is false on a grid
(Remark~N2.2.2, C10 [F]); the proposed sharp version was not sharp when size is
informative (Remark~N2.2.1, C10b [F]); (e) gained a chart and weighted
covariances.
\end{revnote}

\begin{result}[Conjecture N2.2$'$ --- PHASE2c \S4, Remark N2.2.1]\label{res:n22prime}
\status{Conjecture.} The sharp N2 set is
$\{m:\ m(X)\in\E[S(X,T)\mid V]\ P_V\text{-a.s.}\}$ plus the shape restrictions.
A sufficiency argument under the selection-expectation reading (due to the
partial-identification referee, round~2) is recorded as a remark; necessity in
general is not written, and a procedure that uses only $\E[T\mid X]$ does not
attain the set when size is informative. The conjecture stays open.
\end{result}
\begin{plain}
When unit size carries information, conditioning on composition alone can admit
mean rates that the full model excludes (example: $(0.75,0.25)$ admitted,
$(0.5,0.5)$ forced). The same caution applies to the planned M-SHARP series
bound (R9): it is ``tighter than $\ThO$'', not sharp.
\end{plain}
\begin{revnote}(P-B10; I.11.) Added as a conjecture in round~1; the referee's sufficiency remark added in round~2, conjecture kept open by decision.\end{revnote}

\begin{result}[Proposition N2.3, structural local identity --- PHASE2c \S4]\label{res:n23}
\status{Proved} (as a calculus identity). For differentiable $m_{1c},m_{2c}$ and the
structural aggregate $\tau_m(z)=zm_1(z)+(1-z)m_2(z)$: (a)
$\tau'_m=D(x)+\delta_W(x)$; (b)--(c) local signs and local CR in terms of
$\tau'_m$; (d) if both rows are nondecreasing on a neighbourhood of $x$, the
local sign is $\ge0$; monotonicity on the finite support is not enough, and the
local sign everywhere neither implies nor is implied by the aggregate $\delta^*_W$
sign. On a finite support $\tau'$ is not identified, so N2.3 is \emph{not} an
observable restriction; the N2-local rows were dropped (decision U4(a)).
\end{result}
\begin{plain}
EGGHM's local results are a statement about smooth functions, but the data never
identify a derivative at isolated compositions. So ``local'' sign beliefs cannot be
imposed from aggregates without an identified regression family.
\end{plain}
\proofref{Appendix~\ref{app:phase2c}, \S4, Proposition~N2.3 steps 1--3; checks C6, C11, C13, C16}
\begin{revnote}
(P-B11, P-B12, P-B34; I.4.) N2.3 was stated as an observable LP row with an
estimated $\tau'$ (C11 [F]); the smooth witness replaced a kinked one (C13 [I]);
in round~2 (d)'s premise became monotonicity on a neighbourhood (C16 [F]).
\end{revnote}

\begin{result}[Lemma N2.4, person-weighted bins --- PHASE2c \S4]\label{res:n24}
\status{Proved.} With person-weighted bin rates
$\bar m_{1b}=E_b[Xm_1]/\bar x_b$, $\bar m_{2b}=E_b[(1-X)m_2]/(1-\bar x_b)$ and the
model-implied bin mean $\bar\tau_b=E_b\tau$: (i) $\bar m_b\in S(\bar x_b,\bar\tau_b)$;
(ii) the targets are linear in the bin rates; (iii) monotonicity is preserved
across ordered bins. The binned LP is an outer relaxation of global N2.
\end{result}
\begin{plain}
With many compositions, group them into bins, but average rates weighting by
people in each group, not by units; the ordinary bin average breaks the
accounting identity (0.3125 against 0.25 in the referee's example). In samples the
observed bin mean of $T$ estimates $\bar\tau_b$ with error.
\end{plain}
\proofref{Appendix~\ref{app:phase2c}, \S4, Lemma~N2.4 proof; check C12}
\begin{revnote}(P-B13; I.6.) Replaces an invalid [CE-cvx] argument; round~2 distinguished the model-implied bin mean from the observed one.\end{revnote}

\paragraph{What Phase~3 must supply for N2 (PHASE2c \S6).} A weighted polynomial of
\emph{prespecified} degree $q$ (not chosen by cross-validation: post-selection
coverage is about 0.68 in the referee's design, check C12b), conditional on the
specification; then P3.5 applies to the coefficients. (P-B14, I.6.)

\subsection{Results from the package design (PACKAGE\_DESIGN \S\S2.2, 5)}

These results are in Appendix~\ref{app:pkgdesign}. Write
$q_{rs}(x)=(Qx)_r-(Qx)_s$ (identified) and $g_{rs}(x;k)=(Kx)_r-(Kx)_s$.

\begin{result}[G2, JKST plus signs in $2\times2$ --- PACKAGE\_DESIGN \S2.2; register R13]\label{res:g2}
\status{Proved.} Under LCE in $2\times2$ with $\gamma<1$: $\delta_W\ge0$ gives
$[L_J,\min(U_J,D_{ER})]$; $\delta_B\ge0$ gives $[\max(L_J,D_{NM}),U_J]$; CR gives
$[L_J,U_J]\cap[\min,\max](D_{ER},D_{NM})$, where $[L_J,U_J]$ is the JKST interval
for the gap.
\end{result}
\begin{plain}
Adding a sign belief to JKST cuts its interval at one of two data-computed
numbers. In 2000 random LCE designs with true signs, the within sign tightened
JKST in 83\% of designs (median width ratio 0.70), the between sign in 89\%
(0.67). A sign is redundant exactly when the JKST-implied interval for that
association already excludes zero.
\end{plain}
\proofref{Appendix~\ref{app:pkgdesign}, \S2.2, steps 1--3; [V1], [V2]}

\begin{result}[Proposition P-S, sandwich lemma --- PACKAGE\_DESIGN \S5.1; register R12]\label{res:ps}
\status{Proved.} For any $R$ and one column, with the between-group gap
$O_{rs}(x)=(Hx)_r-(Hx)_s$ and the weighted contextual derivative
$L_{rs}(x)=(x^\top H)_r-(x^\top H)_s$: (a) $O_{rs}+L_{rs}=2q_{rs}=\partial_{e_r-e_s}\tau$;
(b) $O_{rs}=q_{rs}+g_{rs}$, $L_{rs}=q_{rs}-g_{rs}$; (c) $O_{rs}\ge0$ and
$L_{rs}\ge0$ at $x$ iff $|g_{rs}(x;k)|\le q_{rs}(x)$, possible only if
$q_{rs}(x)\ge0$ (a data test).
\end{result}
\begin{plain}
An ordering belief (``group $r$ at least as supportive as $s$'') and a local
contextual belief restrict the same unknown from opposite sides, so together
they confine it to an interval whose half-width is an identified slope.
\end{plain}
\proofref{Appendix~\ref{app:pkgdesign}, \S5.1, steps 1--4; [V4]}

\begin{result}[Proposition P-R, rank lemma --- PACKAGE\_DESIGN \S5.2; register R5]\label{res:pr}
(a) One pair restricts an $(R-1)$-dimensional subspace of $\Kspace_c$: \status{proved}.
(b) A star tree of pairs reaches all $\binom R2$ directions: \status{proved}.
(c) Any spanning tree does: \status{checked} for chains, $R=3,4,5$ [V7];
general proof open (R5). (d) M-ORD $\cap$ M-LOC on a star tree bounds $K_c$ without
DD rows: follows from (b) and P-S.
\end{result}
\proofref{Appendix~\ref{app:pkgdesign}, \S5.2}

\begin{result}[Proposition P-N and Conjecture C1, nesting in $2\times2$ under M0 --- PACKAGE\_DESIGN \S5.3; register R1]\label{res:pn}
\status{Proved} for (a)--(b); (c) \status{conjecture}. Under pooled M0: (a) signed
slopes imply weighted-local on the hull; (b) $\delta_W=(a+b)(w_1+d_1x^*)$, so
weighted-local implies $\delta_W\ge0$ whenever $x^*=a/(a+b)$ lies in the
composition hull; (c) Conjecture C1: $x^*$ always lies in the hull (no violation
in 2000 random designs and an adversarial search of 50{,}000).
\end{result}
\begin{revnote}
(P-B7.) The chain ``signed slopes $\subset$ weighted-local $\subset\delta_W\ge0$''
is scoped to M0; it fails for pooled M1 (check I2 of the plan).
\end{revnote}
